\chapter{La recette}\label{ch:recette}

\begin{objectifs}[Objectifs --- étape 5, \texttt{forge-recette}]
\begin{itemize}
  \item dériver le plan de recette des scénarios Gherkin ;
  \item vérifier les critères d'entrée avant de commencer ;
  \item classer les anomalies et gérer leurs délais ;
  \item préparer un procès-verbal que le client, et lui seul, signe.
\end{itemize}
\end{objectifs}

\begin{situation}[Sur le terrain : \sika{}]
Trois trésorières de groupe testent \sika{} sur leur propre téléphone, avec des données anonymisées. L'une d'elles
constate qu'une cotisation confirmée par l'opérateur n'apparaît pas dans le solde du pot. Est-ce grave ?
Réponse du tableau de classification : oui, \emph{bloquante}, car elle empêche le processus métier sans contournement.
\end{situation}

\section{Le concept : la recette est celle du client}
\begin{concept}[Recette]
La phase où le client vérifie, scénario par scénario, que le produit fait ce qui a été validé. Elle ne se confond pas avec
les tests de l'équipe : elle se tient dans un environnement de préproduction, avec des testeurs métier, et se clôt par un
\textbf{procès-verbal (PV)} que le client signe.
\end{concept}

\section{Pas à pas avec \texttt{forge-recette}}
\begin{enumerate}
  \item Le skill lit le plan de recette et chaque \cmd{.feature} des specs du lot.
  \item \textbf{Un scénario Gherkin = une ligne} de la table des scénarios (identifiant, spec, scénario, résultat attendu). Les colonnes \emph{Obtenu}, \emph{Statut}, \emph{Testeur} et \emph{Date} restent vides pour les testeurs.
  \item Il vérifie les \emph{critères d'entrée}.
  \item Il suit les anomalies ; il ne reclasse jamais une anomalie sans l'avis du responsable recette client.
  \item Il prépare le PV ; il ne déclare jamais la recette réussie.
\end{enumerate}

\begin{concept}[Critères d'entrée]
Trois conditions avant de commencer : les specs du lot sont \emph{Validée} ; les scénarios automatisés sont verts en CI ;
un jeu de données anonymisé est chargé. Sinon, la recette ne démarre pas.
\end{concept}

\section{Anomalies et critères de sortie}
\noindent\begin{tabularx}{\linewidth}{llX}
\toprule
\textbf{Niveau} & \textbf{Délai} & \textbf{Définition} \\ \midrule
Bloquante & 48\,h ouvrées & empêche un processus métier sans contournement \\
Majeure & 5 jours ouvrés & processus dégradé, contournement possible \\
Mineure & lot suivant & gêne ergonomique ou cosmétique \\
\bottomrule
\end{tabularx}
\textbf{Critères de sortie :} aucune anomalie bloquante, un nombre borné de majeures avec plan de correction, exigences non fonctionnelles critiques vérifiées.

\subsection*{Un extrait du plan de \sika{}}
\noindent\begin{tabularx}{\linewidth}{lllXl}
\toprule
\textbf{Id} & \textbf{Spec} & \textbf{Scénario} & \textbf{Résultat attendu} & \textbf{Statut} \\ \midrule
SC-01 & SPEC-001 & Cotisation au montant du groupe & cotisation initiée auprès de l'opérateur & OK \\
SC-02 & SPEC-001 & Confirmation par l'opérateur & solde du pot $+$\FCFA{5\,000}, SMS reçu & \textcolor{brick}{KO} \\
SC-03 & SPEC-001 & Seconde cotisation dans le tour & refusée : \q{Déjà cotisé} & OK \\
\bottomrule
\end{tabularx}

\begin{piege}[Piège : corriger à la hâte]
SC-02 échoue. Remplacer l'affichage du solde ferait disparaître le symptôme sans traiter la cause : ici, un webhook de l'opérateur
perdu. Le skill impose la même discipline qu'ailleurs : \emph{reproduire}, poser des hypothèses sur la cause racine, corriger
le minimum via \skill{forge-implemente}, rejouer le scénario. La cause racine retenue : aucune réconciliation périodique des
statuts auprès de l'opérateur. La correction : un rattrapage qui interroge l'opérateur sur les cotisations \q{initiées}.
\end{piege}

\section{Le procès-verbal}
\begin{code}
\begin{Verbatim}
Lot : 1 · Date : [JJ/MM/AAAA] · Décision : [ ] Acceptée [ ] Acceptée avec réserves [ ] Refusée
Réserves :
Signatures : Client ____________ Simetriik Solutions ____________
\end{Verbatim}
\end{code}
Le PV est la porte d'entrée de la production (chapitre~\ref{ch:prod}). Aucune tâche, aucun assistant ne signe à la place du client.

\begin{vousjouez}[À vous de jouer]
Prenez les scénarios de la spec de votre projet et transformez-les en lignes de plan de recette. Qui, chez votre client, jouera le
rôle de responsable recette ? Et quelle anomalie classeriez-vous \emph{bloquante} dans votre contexte ? Écrivez-la.
\end{vousjouez}

\begin{retenir}[À retenir]
\begin{itemize}
  \item Un scénario Gherkin = une ligne du plan de recette.
  \item Critères d'entrée avant de démarrer ; critères de sortie avant de conclure.
  \item Le client classe les anomalies ; le client signe le PV.
\end{itemize}
\end{retenir}
\source{skills/forge-recette/SKILL.md, templates/docs/05-recette}

\begin{acquis}
  \item Citez les trois critères d'entrée d'une recette.
  \item Qui peut reclasser une anomalie ?
  \item Que fait-on d'une anomalie bloquante découverte en recette ?
\end{acquis}
